\chapter{Genital Thrush in Males}
\index{Genital Thrush in Males}

\section*{Definition and Overview}
Genital thrush in males, also known as balanitis due to candida, is a common fungal infection of the penis caused by an overgrowth of the yeast \textit{Candida}, most often \textit{Candida albicans}. Candida organisms live naturally on the skin and in the body in small numbers, but certain conditions allow them to multiply and cause symptoms. In men, thrush typically affects the head of the penis (glans) and the foreskin, causing redness, itching, and sometimes a patchy rash or white discharge. It is not a sexually transmitted infection in the usual sense, although sexual contact can trigger or spread it, and it is readily treatable with antifungal creams or oral medicines.

\section*{Epidemiology}
Genital thrush is less common in men than in women, but it is still a frequent reason for presentation to GPs and sexual health clinics. It occurs most often in men who are uncircumcised, because the warm, moist environment under the foreskin favours yeast growth. Risk is increased in men with diabetes, particularly poorly controlled diabetes, in those taking antibiotics, and in people with weakened immune systems. Thrush is more common in men whose female partners have recurrent vaginal thrush, although most cases are not acquired sexually. Recurrence is common when underlying risk factors are not addressed.

\section*{Aetiology and Pathophysiology}
Genital thrush occurs when candida, which normally lives harmlessly on the skin, multiplies beyond normal levels and invades the superficial layers of the skin of the penis.

\begin{itemize}
    \item \textbf{Normal flora:} candida is part of the normal microbial community of the skin, mouth, and gut
    \item \textbf{Antibiotics:} kill competing bacteria, allowing yeast to overgrow
    \item \textbf{Diabetes:} glucose in the urine and on the skin feeds yeast and impairs immunity
    \item \textbf{Immunosuppression:} HIV, chemotherapy, and corticosteroid treatment reduce the body's ability to control candida
    \item \textbf{Skin conditions:} damage to the skin barrier from eczema, trauma, or poor hygiene favours infection
    \item \textbf{Poor hygiene or moisture:} infrequent washing, tight clothing, or not drying properly after washing
\end{itemize}

\section*{Clinical Features}
Symptoms of genital thrush in males usually develop over a few days and are confined to the penis.

\begin{itemize}
    \item Redness, soreness, and swelling of the glans (head) of the penis
    \item Itching or burning, especially after urination or sexual activity
    \item A patchy rash on the glans or foreskin, sometimes with small white spots
    \item A thick, white, curd-like discharge beneath the foreskin
    \item Difficulty retracting the foreskin (phimosis) due to swelling
    \item Discomfort or pain during sex
    \item An unpleasant odour in some cases
\end{itemize}

\section*{Differential Diagnosis}
Other conditions can cause similar symptoms on the penis, and distinguishing them is important for treatment.

\begin{itemize}
    \item \textbf{Balanitis from other causes:} bacterial infection, irritant dermatitis from soaps or spermicides, and allergic reactions to latex condoms
    \item \textbf{Lichen planus and lichen sclerosus:} inflammatory skin conditions of the penis
    \item \textbf{Psoriasis:} can affect the glans with a distinctive appearance
    \item \textbf{Genital herpes:} causes painful blisters and ulcers
    \item \textbf{Syphilis:} a painless ulcer (chancre) in early infection
    \item \textbf{Genital warts:} raised, wart-like lesions rather than a diffuse rash
    \item \textbf{Reiter syndrome:} urethritis with joint and eye inflammation
\end{itemize}

\section*{When to Seek Medical Care}
See a doctor if symptoms persist beyond a few days of self-care, if they are severe or worsening, or if you have recurrent episodes. Men with diabetes, or those who are immunocompromised, should seek medical advice early.

\textbf{Contact your local emergency services for:}
\begin{itemize}
    \item Sudden severe swelling of the penis with inability to retract the foreskin and difficulty passing urine, which may indicate a serious infection or paraphimosis
    \item Fever with severe pain, redness, or discharge from the penis, which may indicate a spreading bacterial infection
    \item Inability to pass urine at all
\end{itemize}

\section*{Diagnosis and Investigations}
Genital thrush is often diagnosed on the basis of the appearance and symptoms, but tests may be used when the diagnosis is unclear or recurrent.

\begin{itemize}
    \item \textbf{Clinical examination:} inspection of the penis and foreskin
    \item \textbf{Microscopy:} a swab or scraping examined under a microscope can show yeast cells and hyphae
    \item \textbf{Culture:} a swab sent to the laboratory confirms the species and, in resistant cases, its sensitivity to antifungal medicines
    \item \textbf{Blood glucose testing:} recommended if thrush is recurrent, as it may be the first sign of diabetes
    \item \textbf{STI screening:} considered when there is any doubt about the diagnosis or risk of other infections
\end{itemize}

\section*{Management}
Genital thrush in males is usually straightforward to treat, and most men improve quickly.

\begin{itemize}
    \item \textbf{Antifungal creams:} clotrimazole or miconazole cream applied to the glans and foreskin for one to two weeks
    \item \textbf{Oral antifungals:} a single dose of fluconazole is an alternative for more severe or recurrent infection
    \item \textbf{Good hygiene:} washing daily with plain warm water, drying thoroughly, and avoiding soaps and fragranced products
    \item \textbf{Addressing risk factors:} optimising blood glucose control in diabetes and reviewing antibiotic use
    \item \textbf{Treating partners:} women with symptomatic vaginal thrush may benefit from treatment, although routine treatment of asymptomatic partners is not required
    \item \textbf{Recurrence prevention:} if infections return frequently, longer or repeated courses of antifungal treatment may be prescribed
\end{itemize}

\section*{Complications}
\begin{itemize}
    \item Recurrent or chronic infection that is difficult to eradicate
    \item Phimosis, where the foreskin cannot be retracted due to swelling or scarring
    \item Paraphimosis, where a retracted foreskin becomes trapped behind the glans, requiring urgent treatment
    \item Secondary bacterial infection of irritated skin
    \item Spread to the groin or thighs
    \item In immunocompromised men, invasive candidiasis affecting internal organs, which is rare but serious
\end{itemize}

\section*{Prognosis and Outcomes}
With appropriate treatment, most men with genital thrush improve within a few days and clear completely within one to two weeks. Recurrence is common, especially when risk factors such as diabetes or antibiotic use persist. In such cases, addressing the underlying cause and sometimes using longer courses of antifungal therapy are effective. Thrush does not cause long-term damage to the penis or fertility, and once resolved it leaves no lasting effects.

\section*{Prevention and Self-Care}
\begin{itemize}
    \item Wash the penis daily with plain water and dry thoroughly, including under the foreskin
    \item Avoid harsh soaps, bubble baths, and perfumed products in the genital area
    \item Wear loose-fitting cotton underwear and avoid tight clothing and synthetic fabrics
    \item Change out of wet swimwear promptly
    \item If you have diabetes, maintain good blood glucose control
    \item Use condoms consistently, which reduces the risk of transmitting candida and other infections
    \item Seek early treatment for symptoms rather than waiting for them to resolve
    \item See a doctor if you have recurrent episodes, as an underlying cause may need investigation
\end{itemize}

\section*{Sources and Further Reading}
The following reputable sources provide further information for healthcare students and patients seeking reliable, current advice about genital thrush in males.

\begin{itemize}
    \item Healthdirect Australia. \textit{Balanitis and thrush in men.} https://www.healthdirect.gov.au/
    \item NSW Sexual Health Infolink. \textit{Thrush (candida) in men.} https://www.shil.nsw.gov.au/
    \item Melbourne Sexual Health Centre. \textit{Information about balanitis and thrush.} https://www.mshc.org.au/
    \item NHS. \textit{Thrush in men.} https://www.nhs.uk/
    \item Mayo Clinic. \textit{Balanitis: symptoms and causes.} https://www.mayoclinic.org/
    \item DermNet NZ. \textit{Candidal balanitis.} https://dermnetnz.org/
\end{itemize}
